\section{Exact MDS--CSS transversal logical groups}
\label{sec:exact-logical}

The classification rests on one standard MDS conductor statement.  It
converts every nonzero coordinatewise map between two MDS codes with the same
parameters into an isomorphism and leaves room for at most one projective
map.

The conductor and endomorphism-algebra statements below are valid over every
finite field.  For \(q=p^e\) with \(e>1\) they describe the \(\F_q\)-linear
sector of the trace-symplectic Clifford group.
Section~\ref{subsec:additive-labels} treats arbitrary \(\F_p\)-linear label
maps.  It shows that this sector is the whole fixed-coordinate label group
unless \(C\) is diagonally equivalent to a code over a proper subfield of
\(\F_q\), or \(C^\perp\) is diagonally equivalent to a Galois conjugate
\(\tau C\) that is not diagonally equivalent to \(C\).  It identifies the
group in every case.

\begin{proposition}[Equal-parameter MDS conductors]
\label{prop:diagonal-multiplier-line}
Let \(1\leq k<n\), and let \(E,F\leq\F_q^n\) be linear
\([n,k,n-k+1]_q\) MDS codes.
Every nonzero \(s\in\Cond(E,F)\) has full coordinate support and
\(s\star E=F\).  Consequently
\[
 \dim_{\F_q}\Cond(E,F)\leq1,
 \qquad
 \Cond(E,E)=\F_q(1,\ldots,1).
\]
In particular, \(\Cond(C,C^\perp)\) has dimension zero or one,
and \(C\) is diagonally equivalent to its dual exactly in the
dimension-one case.
When that line exists, any two nonzero witnesses differ by a unique
nonzero scalar, so every ratio \(s_i/s_j\) is intrinsic.
\end{proposition}

\begin{proof}
Let \(s\ne0\) have support \(S\) of size \(w\).  Multiplication by \(s\)
on \(E\) has rank \(\min\{k,w\}\), the rank of puncturing an MDS code to
\(S\).  Its image lies in the subcode of \(F\) supported on \(S\), whose
dimension is \(\max\{0,k+w-n\}\).  Hence
\[
 \min\{k,w\}\leq\max\{0,k+w-n\}.
\]
The right side is positive.  If \(w<k\), the inequality would give
\(w\leq k+w-n\), contrary to \(k<n\).  Thus \(w\geq k\), and the same
inequality gives \(w\geq n\).  Therefore \(s\) has full support.
Multiplication by \(s\) is invertible on the ambient space, and the inclusion
\(s\star E\subseteq F\) is equality by equal dimension.

If two conductor elements were independent, a nonzero linear combination could
be chosen to vanish at one coordinate, contradicting full support.
Hence \(\dim\Cond(E,F)\leq1\).  The all-ones vector belongs to
\(\Cond(E,E)\), proving the self-conductor statement.  Taking
\(F=C^\perp\) gives the remaining assertions.
\end{proof}

The standard conductor identity
\begin{equation}\label{eq:conductor-schur-dual}
 \Cond(E,F)=(E\star F^\perp)^\perp
\end{equation}
follows directly from the Euclidean inner product.  Together with the product
Singleton bound, it gives another proof of the dimension statement
\cite{CouvreurMarquezCorbellaPellikaan2017,MirandolaZemor2015}.  We retained
the elementary proof to make the precise MDS input visible.  Specializing
\eqref{eq:conductor-schur-dual} gives
\begin{equation}\label{eq:dual-conductor-schur-square}
 \Cond(C,C^\perp)=(C^{\star2})^\perp.
\end{equation}
Thus the code-to-dual test has only two outcomes: no conductor, or a unique
full-support conductor line.  Equivalently, the Schur square has codimension
zero or one.

Fix an input coordinate \(j\).  Under the Choi correspondence, a physical
product unitary implements a logical unitary \(L\) exactly when adjoining the
input factor \(L^{-T}\) gives a \(2m\)-party product symmetry of
\(\lvert\Psi_C\rangle\).  We therefore classify the local symplectic blocks
that preserve the state label space \(L_C\).

Write a fixed-coordinate product-Clifford action on Pauli labels as
\[
 F_i=\begin{pmatrix}\alpha_i&\beta_i\\
                     \gamma_i&\delta_i\end{pmatrix},
\]
and let \(D_\alpha,D_\beta,D_\gamma,D_\delta\) be the diagonal
matrices with the corresponding entries.  Preservation of
\(L_C=C_X\oplus C^\perp_Z\) is equivalent to
\begin{equation}\label{eq:block-preservation}
\begin{array}{ll}
D_\alpha C\subseteq C,&D_\gamma C\subseteq C^\perp,\\
D_\beta C^\perp\subseteq C,&D_\delta C^\perp\subseteq C^\perp.
\end{array}
\end{equation}
These are respectively the \(X\mathbin\to X\), \(X\mathbin\to Z\),
\(Z\mathbin\to X\), and \(Z\mathbin\to Z\) preservation conditions.

To collect the four standard conductor spaces in one object, define the
\emph{algebra of coordinatewise endomorphisms}
\begin{equation}\label{eq:css-endomorphism-algebra}
 \mathcal A_C=
 \begin{pmatrix}
  \Cond(C,C)&\Cond(C^\perp,C)\\
  \Cond(C,C^\perp)&\Cond(C^\perp,C^\perp)
 \end{pmatrix}
 \subseteq M_2(\F_q^{2m}),
\end{equation}
with coordinatewise matrix multiplication.

\begin{proposition}[The two coordinatewise endomorphism algebras]
\label{prop:two-endomorphism-algebras}
The algebra \(\mathcal A_C\) is exactly the algebra of sitewise
\(\F_q\)-linear Pauli-label maps preserving
\(L_C=C_X\oplus C_Z^\perp\), and
\[
 \Cond(C,C)=\Cond(C^\perp,C^\perp)=\F_q\mathbf1.
\]
If \(\Cond(C,C^\perp)=0\), then
\(\mathcal A_C\cong\F_q\times\F_q\).  Otherwise, after choosing
\(s\star C=C^\perp\), there is an algebra isomorphism
\[
 M_2(\F_q)\longrightarrow\mathcal A_C,
 \qquad
 \begin{pmatrix}a&b\\c&d\end{pmatrix}
 \longmapsto
 \begin{pmatrix}a\mathbf1&b s^{-1}\\cs&d\mathbf1\end{pmatrix}.
\]
The coordinatewise determinant-one units are respectively \(T\) and
\(\mathrm{SL}_2(q)\).
\end{proposition}

\begin{proof}
The four entries in \eqref{eq:css-endomorphism-algebra} are precisely the
four inclusions in \eqref{eq:block-preservation}.  They are closed under
coordinatewise matrix multiplication: for example,
\[
 \Cond(C^\perp,C)\star\Cond(C,C^\perp)\subseteq\Cond(C,C),
\]
because multiplication first sends \(C\) into \(C^\perp\) and then back
into \(C\).  The other products are identical.  The conductor identity
\eqref{eq:conductor-schur-dual} gives
\[
 \Cond(C,C)=(C\star C^\perp)^\perp
 =\Cond(C^\perp,C^\perp),
\]
and Proposition~\ref{prop:diagonal-multiplier-line} makes this common algebra
the scalar line.  If \(\Cond(C,C^\perp)=\F_qs\) is nonzero, then \(s\) has
full support and \(\Cond(C^\perp,C)=\F_qs^{-1}\); the converse follows in
the same way.  Thus the two off-diagonal conductors either both vanish or are
these two inverse lines.  This gives the two displayed algebra types.  In the
matrix-algebra case the determinant at every coordinate is \(ad-bc\); in the
split case it is \(ad\).  Imposing determinant one proves the final statement.
\end{proof}

The proposition packages the block calculation into the split-algebra versus
matrix-algebra dichotomy: the MDS property leaves no intermediate
coordinatewise linear symmetry algebra.  We now use its two cases to prove the
logical-group theorem.

\begin{proof}[Proof of
Theorem~\ref{thm:diagonal-isodual-transversal-group}]

\medskip\noindent\emph{Zero-conductor case.}
Proposition~\ref{prop:two-endomorphism-algebras} identifies every
label-preserving symplectic action with an element of \(T\), realized by the
same block \(\diag(a,a^{-1})\) at every coordinate.  Conversely, these blocks
preserve \(L_C\).  Logical Paulis have physical Pauli representatives, so this
branch gives \(\F_q^2\rtimes T\).

\medskip\noindent\emph{Nonzero-conductor case.}
Now suppose
\begin{equation}\label{eq:diagonal-duality-witness}
 S=\diag(s_1,\ldots,s_{2m}),\qquad SC=C^\perp .
\end{equation}
The projective vector \([s_1:\cdots:s_{2m}]\), and hence every propagation
ratio \(s_i/s_j\), depends only on \(C\).

For \(\lambda,\mu\in\F_q\), put
\begin{equation}\label{eq:duality-shears}
 F_i^-(\lambda)=
 \begin{pmatrix}1&0\\ \lambda s_i&1\end{pmatrix},
 \qquad
 F_i^+(\mu)=
 \begin{pmatrix}1&\mu s_i^{-1}\\0&1\end{pmatrix}.
\end{equation}
On the CSS label space
\(L_C=(C\oplus0)+(0\oplus C^\perp)\), the product of the lower blocks
sends
\[
 (c,h)\longmapsto(c,h+\lambda Sc),
\]
and the product of the upper blocks sends
\[
 (c,h)\longmapsto(c+\mu S^{-1}h,h).
\]
Thus both products preserve \(L_C\) and lift to product-Clifford tensor
symmetries: Lemma~\ref{lem:pauli-phase-correction} supplies the required
product-Pauli corrections.  Thus label-space preservation is sufficient
here; no coherent choice of Clifford phases is being assumed.
At any chosen input party \(j\), their blocks run through all lower and
upper elementary unipotents because \(s_j\ne0\).  Together these generate
\(\mathrm{SL}_2(q)\).  Under the Choi correspondence, an input unitary
\(U\) implements the logical unitary \(U^{-T}=\overline{U}\).  Complex
conjugation fixes \(X(a)\) and sends \(Z(b)\) to \(Z(-b)\), so an input
block \(A\) induces the logical block \(RAR\) with \(R=\diag(1,-1)\).
This map is an automorphism of \(\mathrm{SL}_2(q)\) preserving the
diagonal torus, so every one-qudit logical symplectic action has a
product-Clifford physical implementation.  Logical Paulis have
physical Pauli representatives, so the full projective one-qudit
Clifford group
\(\F_q^2\rtimes\mathrm{SL}_2(q)\) is transversal.

\medskip\noindent\emph{Exactness.}
The imported Corollary~\ref{cor:transversal-clifford} excludes every product-unitary
logical action outside that group.  Hence the displayed group is exact.
\end{proof}

\AMEFigureThree

For a generalized or extended generalized Reed--Solomon code, the
standard dual formula supplies \eqref{eq:diagonal-duality-witness}, including the extended
evaluation point at length \(q+1\)
\cite[Sections~5.2--5.3]{HuffmanPless2003}.  Thus the GRS tower lies on the
branch diagonally equivalent to its dual, although the proof uses no
evaluation structure.

\begin{proposition}[Site-dependent propagation]
\label{prop:site-dependent-realization}
Fix an input coordinate \(j\), put \(t_i=s_i/s_j\), and write
\[
 A=\begin{pmatrix}a&b\\c&d\end{pmatrix}\in\mathrm{SL}_2(q).
\]
Propagation to coordinate \(i\) is the twisted diagonal embedding
\begin{equation}\label{eq:site-propagation}
 \phi_i(A)=
 \begin{pmatrix}a&t_i^{-1}b\\t_i c&d\end{pmatrix}
 =
 \begin{pmatrix}1&0\\0&t_i\end{pmatrix}
 A
 \begin{pmatrix}1&0\\0&t_i^{-1}\end{pmatrix}.
\end{equation}
\end{proposition}

\begin{proof}
The formula follows from the upper and lower elementary matrices in
\eqref{eq:duality-shears}, which generate \(\mathrm{SL}_2(q)\).
\end{proof}

The formula is the reason for allowing the one-qudit factor to depend on the
coordinate: a common logical matrix generally propagates differently at
different sites.  Appendix~\ref{sec:lift-refinements} asks when fixed local
frame changes can make these actions uniform and records two phase
refinements.  The linear
\(\mathrm{SL}_2(q)\) factor admits coherent Weil lifts, whereas adjoining
logical Pauli translations gives the nonsplit Heisenberg central extension.

\subsection{Prime-power fields: additive label maps}
\label{subsec:additive-labels}

Let \(q=p^e\) and \(\Gamma=\mathrm{Gal}(\F_q/\F_p)\).  A local Clifford acts
on Pauli labels by an arbitrary trace-symplectic \(\F_p\)-linear map
\eqref{eq:projective-clifford-sequence}, and the block equations
\eqref{eq:block-preservation} hold verbatim for \(\F_p\)-linear blocks.  The
question is therefore which coordinatewise \(\F_p\)-linear maps carry one MDS
code into another.  The answer is governed by a subfield of \(\F_q\) read
off from the code, and by at most one class of field automorphisms.

For a linear code \(E\leq\F_q^n\) and \(\tau\in\Gamma\), put
\(\tau E=\{(\tau c_1,\ldots,\tau c_n):c\in E\}\).  Write \(E\sim F\), and
call \(E\) and \(F\) \emph{diagonally equivalent}, when \(F=s\star E\) for
some \(s\in(\F_q^\times)^n\).  The \emph{additive conductor} is
\[
 \Cond_{\F_p}(E,F)=\{\alpha\in\End_{\F_p}(\F_q)^n:\alpha(E)\subseteq F\},
\]
with \(\alpha=(\alpha_1,\ldots,\alpha_n)\) acting coordinatewise.  If \(E\)
is a linear \([n,k,n-k+1]_q\) MDS code with \(1\leq k<n\), every
\(k\)-subset of coordinates is an information set, and the systematic
generator matrix \([I_k\mid A]\) on an information set \(I\) has no zero
entry, because each entry is a \(k\)-by-\(k\) minor up to sign.  Its
\emph{normalized matrix} is
\[
 N_E=\bigl(A_{ij}A_{11}/(A_{i1}A_{1j})\bigr)_{i,j}.
\]
Its entries are cross-ratios of the entries of \(A\).  Let
\(K_E=\F_p(N_E)\) be the \emph{cross-ratio field}.  When
\(k\in\{1,n-1\}\), the matrix \(A\) is a single row or column, so \(N_E\)
consists of ones and \(K_E=\F_p\).

\begin{lemma}[Cross-ratio field]
\label{lem:cross-ratio-field}
Let \(E,F\leq\F_q^n\) be linear \([n,k,n-k+1]_q\) MDS codes with
\(1\leq k<n\), and let \(\tau\in\Gamma\).
\begin{enumerate}
\item On a fixed information set, \(E\sim F\) if and only if \(N_E=N_F\),
  and \(N_{\tau E}=\tau(N_E)\).
\item \(K_E\) is the smallest subfield \(L\leq\F_q\) for which \(E\) is
  diagonally equivalent to the extension of scalars of a code over \(L\).
  In particular \(K_E\) does not depend on the information set.
\item \(\tau E\sim E\) if and only if \(\tau\) fixes \(K_E\) pointwise.
\item \(K_{E^\perp}=K_E\).
\end{enumerate}
\end{lemma}

\begin{proof}
If \(F=s\star E\), then
\(F=\mathrm{row}[\,I_k\mid\diag(s_I)^{-1}A\diag(s_J)\,]\), where \(J\) is
the complement of \(I\).  Diagonal equivalence therefore acts on systematic
matrices by independent row and column scalings, and \(N_E\) is the unique
matrix in the orbit of \(A\) whose first row and first column consist of
ones.  This proves the first assertion of item~1.  The second holds because
Gaussian elimination commutes with \(\tau\).

The code \(\mathrm{row}[I_k\mid N_E]\) is defined over \(K_E\) and is
diagonally equivalent to \(E\).  Conversely, if \(E\sim E_0\otimes\F_q\) with
\(E_0\leq L^n\), then \(N_E=N_{E_0}\) has entries in \(L\).  This proves
item~2.  Item~3 follows from item~1: \(\tau E\sim E\) exactly when
\(\tau(N_E)=N_E\).  For item~4, \((s\star E)^\perp=s^{-1}\star E^\perp\) and
\((E_0\otimes\F_q)^\perp=E_0^\perp\otimes\F_q\), so \(E^\perp\) is
diagonally equivalent to a code over \(K_E\), and \(K_{E^\perp}\subseteq
K_E\).  The reverse inclusion is the same statement for \(E^\perp\).
\end{proof}

For a subfield \(K\leq\F_q\) and \(\sigma\in\mathrm{Gal}(K/\F_p)\), let
\(\End_K^\sigma(\F_q)\) be the space of \(\sigma\)-semilinear \(K\)-maps
\(\nu\) of \(\F_q\), that is, \(\nu(\lambda u)=\sigma(\lambda)\nu(u)\) for
\(\lambda\in K\).  The next proposition shows that every additive conductor
between equal-parameter MDS codes is one such space, applied with the same
\(\nu\) at every coordinate up to a fixed diagonal twist.

\begin{proposition}[Additive MDS conductors]
\label{prop:additive-conductor}
Let \(1\leq k<n\), and let \(E,F\leq\F_q^n\) be linear
\([n,k,n-k+1]_q\) MDS codes.  Then
\begin{equation}\label{eq:galois-sector-decomposition}
 \Cond_{\F_p}(E,F)=\bigoplus_{\tau\in\Gamma}\Cond(\tau E,F)\,\tau,
 \qquad
 \Cond(\tau E,F)\,\tau=\{x\mapsto s\star\tau(x):s\in\Cond(\tau E,F)\},
\end{equation}
and the following hold.
\begin{enumerate}
\item Every nonzero \(\alpha\in\Cond_{\F_p}(E,F)\) has \(\alpha_i\ne0\) for
  every \(i\).
\item \(\Cond_{\F_p}(E,F)\ne0\) exactly when \(\tau E\sim F\) for some
  \(\tau\in\Gamma\).  In that case \(K_E=K_F=:K\), the restriction
  \(\sigma=\tau|_K\) does not depend on the choice of \(\tau\), the nonzero
  summands are exactly those with \(\tau|_K=\sigma\), and
  \[
   \dim_{\F_p}\Cond_{\F_p}(E,F)=e\,[\F_q:K].
  \]
\item In that case, if \(E=E_0\otimes\F_q\) and \(F=F_0\otimes\F_q\) with
  \(E_0,F_0\leq K^n\), then \(F_0=r\star\sigma(E_0)\) for some
  \(r\in(K^\times)^n\), and
  \[
   \Cond_{\F_p}(E,F)=\{x\mapsto r\star\nu(x):\nu\in\End_K^\sigma(\F_q)\},
  \]
  with \(\nu\) applied at every coordinate.  In particular every element is
  \(\sigma\)-semilinear over \(K\), with the same \(\sigma\) at every
  coordinate.
\end{enumerate}
Consequently \(\Cond_{\F_p}(C,C)\cong\End_{K_C}(\F_q)\), which is
\(\F_q\mathbf1\) when \(K_C=\F_q\).  The code-to-dual conductor
\(\Cond_{\F_p}(C,C^\perp)\) is either zero or concentrated on one class
\(\sigma\in\mathrm{Gal}(K_C/\F_p)\), and then \(\sigma^2=1\).
\end{proposition}

\begin{proof}
By Dedekind's independence of characters, every
\(\phi\in\End_{\F_p}(\F_q)\) is uniquely \(\sum_\tau c_\tau\tau\) with
\(c_\tau\in\F_q\); the maps \(c\tau\) are \(\F_q\)-independent, and both
sides have \(\F_p\)-dimension \(e^2\).  Thus every \(\alpha\) is uniquely
\(\sum_\tau s_\tau\tau\) with \(s_\tau\in\F_q^n\).  Let \(a\) generate
\(\F_q\) over \(\F_p\).  Because \(E\) and \(F\) are \(\F_q\)-linear, the
conductor is stable under \(\alpha\mapsto\alpha m_a\) and under
\(\alpha\mapsto m_c\alpha\) for constants \(c\in\F_q\), where \(m_c\) is
multiplication by \(c\).  On the \(\tau\)-component, the operator
\(\alpha\mapsto\alpha m_a-m_{\rho(a)}\alpha\) is multiplication by
\(\tau(a)-\rho(a)\).  The conjugates of \(a\) are distinct, so the product of
these operators over \(\rho\ne\tau\) kills every other component and
rescales the \(\tau\)-component by a nonzero constant.  Hence each
component \(s_\tau\tau\) of a conductor element is a conductor element, and
\(s_\tau\tau\) lies in the conductor exactly when
\(s_\tau\star\tau(E)\subseteq F\).  This proves
\eqref{eq:galois-sector-decomposition}.

Proposition~\ref{prop:diagonal-multiplier-line} applies to the MDS codes
\(\tau E\) and \(F\).  Each summand is zero or a full-support line, and it
is nonzero exactly when \(\tau E\sim F\).  Item~1 follows: \(\alpha_i=0\)
forces \(s_{\tau,i}=0\) for every \(\tau\), by uniqueness of the expansion,
and hence \(s_\tau=0\) for every \(\tau\).  If \(\tau E\sim F\), then
\(\tau'E\sim F\) if and only if \(\tau^{-1}\tau'E\sim E\), that is, if and
only if \(\tau'|_{K_E}=\tau|_{K_E}\), by Lemma~\ref{lem:cross-ratio-field}.
The nonzero summands therefore form one coset of
\(\mathrm{Gal}(\F_q/K_E)\), which has \([\F_q:K_E]\) elements, each
contributing a line of \(\F_p\)-dimension \(e\).  Also
\(K_F=\F_p(\tau N_E)=\tau(K_E)\), and subfields of \(\F_q\) of equal order
coincide, so \(K_F=K_E\).  This proves item~2.

For item~3, \(E_0\) and \(F_0\) are diagonally equivalent over \(K\) to
\(\mathrm{row}[I_k\mid N_E]\) and \(\mathrm{row}[I_k\mid N_F]\), by the
scalings in the proof of Lemma~\ref{lem:cross-ratio-field}, and
\(N_F=\tau N_E=\sigma N_E\).  Hence \(F_0=r\star\sigma(E_0)\) with
\(r\in(K^\times)^n\).  For every \(\tau'\) with \(\tau'|_K=\sigma\),
\(\tau'E=\sigma(E_0)\otimes\F_q\) and
\(\Cond(\tau'E,F)=\F_q r\).  The conductor is therefore the set of maps
\(r\star\nu\) with \(\nu\) in the span of the maps \(c\tau'\), and this span
is \(\End_K^\sigma(\F_q)\): it is contained there, and both have
\(\F_p\)-dimension \(e[\F_q:K]\).

For \(E=F=C\), take \(\sigma=1\).  For the dual statement,
\(C^\perp=s\star\tau C\) gives
\(C=s^{-1}\star\tau(C^\perp)=s^{-1}\tau(s)\star\tau^2C\), so \(\tau^2\)
fixes \(K_C\) by Lemma~\ref{lem:cross-ratio-field}.
\end{proof}

For invertible coordinate maps, semilinearity over the field generated by
the entries of a standard-form generator matrix is the mechanism behind
Adriaensen and Ball's characterization of additive MDS codes equivalent to
linear codes \cite[Lemmas~5.2 and~5.6]{AdriaensenBall2023}, where coordinate
maps are written as linearised polynomials, that is, expanded in Galois
sectors.  The proposition applies that expansion to the whole conductor,
including its noninvertible elements, which occur whenever \(K\ne\F_q\), and
computes its exact dimension.

Proposition~\ref{prop:additive-conductor} sorts a code \(C\) into three
types: the additive code-to-dual conductor is zero; it is concentrated on
\(\sigma=1\), so that \(C\sim C^\perp\) as in
Proposition~\ref{prop:two-endomorphism-algebras}; or it is concentrated on
the involution \(\sigma\) of \(K_C\), so that \(C^\perp\sim\tau C\) for an
automorphism \(\tau\) with \(\tau C\not\sim C\).  We call the three types
\emph{zero}, \emph{Euclidean}, and \emph{Hermitian}.  In the Hermitian type,
\([K_C:\F_p]\) is even, and \(K_0\) denotes the fixed field of \(\sigma\),
so \(|K_C|=|K_0|^2\).

Write \(G^{\mathrm{add}}_C\) for the group of sitewise \(\F_p\)-linear
Pauli-label maps that preserve \(L_C=C_X\oplus C^\perp_Z\) and the
trace-symplectic form
\(\omega((x,z),(x',z'))=\Tr_{\F_q/\F_p}(xz'-x'z)\) at every coordinate.

\begin{theorem}[Additive fixed-coordinate label groups]
\label{thm:additive-label-groups}
Let \(q=p^e\), let \(C\leq\F_q^{2m}\) be a linear \([2m,m,m+1]_q\) MDS code,
and put \(K=K_C\) and \(d=[\F_q:K]\).  Projection of
\(G^{\mathrm{add}}_C\) to any one coordinate is injective, and
\[
 G^{\mathrm{add}}_C\cong
 \begin{cases}
  \mathrm{GL}_d(K), & C\text{ of zero type},\\
  \mathrm{Sp}_{2d}(K), & C\text{ of Euclidean type},\\
  \mathrm{GU}_{2d}(K_0), & C\text{ of Hermitian type}.
 \end{cases}
\]
In the Euclidean and Hermitian types the sitewise \(\F_p\)-linear maps
preserving \(L_C\) form an algebra isomorphic to \(M_{2d}(K)\), and
\(G^{\mathrm{add}}_C\) is the isometry group of a nondegenerate
\(\sigma\)-skew-Hermitian form on \(K^{2d}\).  If \(K=\F_q\), the group
\(G^{\mathrm{add}}_C\) equals the \(\F_q\)-linear group \(T\) or
\(\mathrm{SL}_2(q)\) of Proposition~\ref{prop:two-endomorphism-algebras} in
the first two types.  In the Hermitian type it is
\(\mathrm{GU}_2(\sqrt q)\), and exactly \(2(q-1)\) of its elements are
semilinear for an automorphism of \(\F_q\).
\end{theorem}

The proof turns the label algebra into the endomorphism algebra of one
\(2d\)-dimensional \(K\)-space.  In the Hermitian type, the \(K\)-structure
on the \(Z\)-labels is twisted by \(\sigma\); this twist is what converts
the trace-symplectic form into a skew-Hermitian one.

\begin{proof}
\emph{Reduction.}  Choose \(C=D(C_0\otimes\F_q)\) with \(D\) invertible
diagonal and \(C_0\leq K^{2m}\), by Lemma~\ref{lem:cross-ratio-field}.  The
sitewise map \((x,z)\mapsto(D^{-1}x,Dz)\) preserves \(\omega\) at every
coordinate and sends \(L_C\) to \(L_{C_0\otimes\F_q}\), because
\((DE)^\perp=D^{-1}E^\perp\).  Conjugation by it identifies the two label
groups, so we may assume \(C=C_0\otimes\F_q\).  Then
\(C^\perp=C_0^\perp\otimes\F_q\).  The four blocks of a label map are
additive conductors, and Proposition~\ref{prop:additive-conductor}, item~3,
describes each of them.  The diagonal blocks are maps \(\mu\) and \(\mu'\)
in \(\End_K(\F_q)\), applied at every coordinate.  Injectivity of the
projection follows from item~1 of that proposition.

\emph{Zero type.}  Both off-diagonal blocks vanish: a nonzero conductor
\(C^\perp\to C\) would give \(\tau C^\perp\sim C\), and dualizing would give
\(\tau C\sim C^\perp\).  Write \(T_K=\Tr_{\F_q/K}\).  A pair
\((\mu,\mu')\) preserves \(\omega\) exactly when
\(\Tr_{K/\F_p}T_K(\mu(x)\mu'(z))=\Tr_{K/\F_p}T_K(xz)\) for all \(x,z\).  Both
sides are \(\Tr_{K/\F_p}\) of a \(K\)-bilinear form, and replacing \(x\) by
\(\lambda x\) with \(\lambda\in K\) shows that this holds exactly when the
two \(K\)-bilinear forms agree.  Thus \(\mu'\) is the inverse adjoint of
\(\mu\) for the nondegenerate pairing \(T_K(xz)\), and
\(\mu\mapsto(\mu,\mu')\) identifies \(G^{\mathrm{add}}_C\) with
\(\mathrm{GL}_K(\F_q)\cong\mathrm{GL}_d(K)\).

\emph{Euclidean and Hermitian types.}  Let \(\sigma\) be the class of the
code-to-dual conductor, and let \(K_0\) be its fixed field, with \(K_0=K\)
when \(\sigma=1\).  By Proposition~\ref{prop:additive-conductor},
\(C_0^\perp=r\star\sigma(C_0)\) with \(r\in(K^\times)^{2m}\).  Dualizing
gives \(C_0=r^{-1}\star\sigma(C_0^\perp)\).  Applying \(\sigma\), and using
\(\sigma^2=1\), gives \(C_0^\perp=\sigma(r)\star\sigma(C_0)\).  Hence
\(\sigma(r)=\lambda r\) for a scalar \(\lambda\in K^\times\), by
Proposition~\ref{prop:diagonal-multiplier-line} over \(K\).  Applying
\(\sigma\) again gives \(\lambda\sigma(\lambda)=1\).  By Hilbert's
Theorem~90 for \(K/K_0\) (trivial when \(\sigma=1\)),
\(\lambda=\kappa/\sigma(\kappa)\) for some
\(\kappa\in K^\times\), and \(\kappa r\) is fixed by \(\sigma\).  Replacing
\(r\) by \(\kappa r\), we may assume \(r_i\in K_0^\times\) for every \(i\).
Proposition~\ref{prop:additive-conductor} now gives the off-diagonal blocks
\(x\mapsto r\star\nu'(x)\) and \(z\mapsto r^{-1}\star\nu(z)\), with
\(\nu,\nu'\in\End_K^\sigma(\F_q)\) applied at every coordinate.

Let \(W=\F_q\oplus\F_q\) with the \(K\)-structure
\(\lambda\cdot(x,z)=(\lambda x,\sigma(\lambda)z)\).  A map
\(h=\begin{psmallmatrix}\mu&\nu\\\nu'&\mu'\end{psmallmatrix}\) of \(W\) is
\(K\)-linear exactly when \(\mu,\mu'\) are \(K\)-linear and \(\nu,\nu'\) are
\(\sigma\)-semilinear.  Put \(\Delta_i(x,z)=(x,r_iz)\).  Because
\(r_i\in K_0\), the label map at coordinate \(i\) is
\(\Delta_ih\Delta_i^{-1}\) for one \(h\) independent of \(i\), and
\(h\) ranges over all of \(\End_K(W)\cong M_{2d}(K)\).  The map from label
maps to \(h\) respects composition.  Define
\[
 H((x,z),(x',z'))=T_K(xz')-\sigma\bigl(T_K(x'z)\bigr).
\]
It is \(K\)-linear in the first variable and \(\sigma\)-semilinear in the
second, it satisfies \(H(w',w)=-\sigma(H(w,w'))\), and it is
nondegenerate because \(T_K(xz)\) is.  Since
\(r_i\in K_0\) and \(\Tr_{K/\F_p}\circ\sigma=\Tr_{K/\F_p}\),
\[
 \omega(\Delta_iw,\Delta_iw')=\Tr_{K/\F_p}\bigl(r_iH(w,w')\bigr).
\]
For \(K\)-linear \(h\), the difference \(H(hw,hw')-H(w,w')\) is again
sesquilinear, and replacing \(w\) by \(\lambda\cdot w\) shows that
\(h\) preserves \(\Tr_{K/\F_p}(r_iH)\) exactly when it preserves \(H\).  The
condition is the same at every coordinate.  Hence \(G^{\mathrm{add}}_C\) is
the isometry group of \(H\) on the \(2d\)-dimensional \(K\)-space \(W\).

If \(\sigma=1\), then \(H\) is alternating, since \(H(w,w)=0\), and its
isometry group is \(\mathrm{Sp}_{2d}(K)\).  If \(\sigma\ne1\), choose
\(\varepsilon\in K^\times\) with \(\sigma(\varepsilon)=-\varepsilon\)
(\(\varepsilon=1\) in characteristic two).  Then \(\varepsilon H\) is a
nondegenerate Hermitian form with the same isometry group.  Because
\(\sigma\ne1\), the trace \(\Tr_{K/K_0}\) is not identically zero, and for a
Hermitian form \(h\) polarization gives
\[
 h(u+\lambda v,u+\lambda v)-h(u,u)-N_{K/K_0}(\lambda)h(v,v)
 =\Tr_{K/K_0}\bigl(\lambda h(v,u)\bigr).
\]
Hence a Hermitian
form that is not identically zero has an anisotropic vector, so it is
diagonalizable.  Its diagonal entries lie in \(K_0^\times\), and the
norm \(K^\times\to K_0^\times\) is surjective.  Every nondegenerate
Hermitian form on \(K^{2d}\) is therefore isometric to the standard one, and
its isometry group is \(\mathrm{GU}_{2d}(K_0)\).

If \(K=\F_q\), then \(d=1\), every block is a multiple of one field
automorphism, and in the first two types \(\sigma=1\), so every label map is
\(\F_q\)-linear.  In the Hermitian type, with the basis
\((1,0),(0,1)\) of \(W\), the matrix of \(h\) is
\(\begin{psmallmatrix}a&b\\\sigma(c)&\sigma(d)\end{psmallmatrix}\) when
\(\mu,\nu,\nu',\mu'\) are multiplication by \(a\), \(b\sigma\), \(c\sigma\),
and \(d\), respectively.  An element is semilinear for one automorphism of
\(\F_q\) exactly when it is diagonal or antidiagonal, by uniqueness of the
Galois expansion.  The isometry conditions reduce to \(ad=1\) and to
\(bc=-1\), respectively, which give \(2(q-1)\) elements.
\end{proof}

Over a prime field, \(K=\F_q\) and \(\sigma=1\), so only the first two
types occur.  The theorem then recovers
Proposition~\ref{prop:two-endomorphism-algebras}.

\begin{corollary}[Transversal groups over prime-power fields]
\label{cor:prime-power-transversal-group}
Let \(q\) be a prime power, let \(m\geq2\), let \(C\leq\F_q^{2m}\) be a
linear \([2m,m,m+1]_q\) MDS code, and fix an input coordinate \(j\).  Let
\(\pi_j\colon G^{\mathrm{add}}_C\to\mathrm{Sp}_{2e}(\F_p)\) be projection
to coordinate \(j\), and let \(R(x,z)=(x,-z)\).  The projective group of
logical unitaries with fixed-coordinate transversal implementations is the
full preimage of the subgroup \(R\,\pi_j(G^{\mathrm{add}}_C)\,R\) under
the label map \(\operatorname{PCliff}(q)\to\mathrm{Sp}_{2e}(\F_p)\) of
\eqref{eq:projective-clifford-sequence}.  This subgroup is isomorphic to
\(G^{\mathrm{add}}_C\), hence to \(\mathrm{GL}_d(K_C)\),
\(\mathrm{Sp}_{2d}(K_C)\), or \(\mathrm{GU}_{2d}(K_0)\) according to the
type of \(C\).  The kernel of the logical group's label map is the group
\(\F_q^2\) of logical Pauli translations.
\end{corollary}

\begin{proof}
\emph{Forward inclusion.}  The imported
Corollary~\ref{cor:transversal-clifford} holds for every prime power and
makes every factor of a product implementation Clifford.  Its label map is
therefore trace-symplectic and \(\F_p\)-linear at every coordinate, and it
preserves \(L_C\), so it lies in \(G^{\mathrm{add}}_C\).

\emph{Lifts.}  Conversely, every element of \(G^{\mathrm{add}}_C\) has a
Clifford lift at each coordinate by
\eqref{eq:projective-clifford-sequence}, which holds for \(p=2\) as well, and
Lemma~\ref{lem:pauli-phase-correction} corrects the product by a product
Pauli to a symmetry of \(\lvert\Psi_C\rangle\).  Logical Paulis have
physical Pauli representatives, so the logical group contains the kernel
\(\F_q^2\) and is the full preimage of its image.

\emph{Choi twist.}  As in the proof of
Theorem~\ref{thm:diagonal-isodual-transversal-group}, an input factor \(U\)
implements the logical unitary \(\overline U\), and complex conjugation
sends \(X(a)\) to \(X(a)\) and \(Z(b)\) to \(Z(-b)\).  The logical label map
is therefore \(R\,\pi_j(g)\,R\) for the label symmetry \(g\), and \(\pi_j\)
is injective by Theorem~\ref{thm:additive-label-groups}.  In characteristic
two the additive character is real-valued, so complex conjugation fixes
every Pauli, \(R\) is the identity, and the logical label map is the input
block itself.  For odd \(p\) the preimage splits by the Weil
representation; for \(p=2\) it need not, so no semidirect product is
asserted.
\end{proof}

For \(q\) an odd prime the corollary is
Theorem~\ref{thm:diagonal-isodual-transversal-group}.  The
hexacode \([6,3,4]_4\) is of Hermitian type with \(K=\F_4\): its
fixed-coordinate label group is \(\mathrm{GU}_2(2)\), of order \(18\), whereas
the \(\F_4\)-linear group is the torus \(T\) of order \(3\).  Of the \(18\)
label symmetries, \(12\) are semilinear for no automorphism of \(\F_4\).

The hexacode is the smallest Hermitian example.  Hermitian type needs
\([K_C:\F_p]\) even, so \(q\geq4\).  For \(m=2\), every \([4,2,3]_q\) code is
generalized Reed--Solomon, hence diagonally equivalent to its dual and of
Euclidean type.  For \(m=3\) and \(q=4\), the columns of a generator matrix
are six points of \(\mathrm{PG}(2,4)\), no three collinear.  Any five of
them form a conic, and every line through the sixth meets that conic in at
most one point, so it is tangent and the sixth point is the nucleus.  Hence
every \([6,3,4]_4\) code is monomially equivalent to the hexacode.
